% VI25_CHALLENGE_SOLUTIONS_REFINEMENT_V1
\subsection*{Challenge 1: Complete degeneration audit}
\begin{solution}
Consider
\[
X=V\bigl(y(y-tx-1)\bigr)\subset\mathbb A^2_{x,y}\times\mathbb A^1_t.
\]
For \(c\in k\), the fiber is
\[
X_c=\Spec k[x,y]/\bigl(y(y-cx-1)\bigr).
\]
The two factors
\[
y,\qquad y-cx-1
\]
are distinct irreducible linear polynomials for every \(c\).  Hence their product is squarefree.  In the UFD \(k[x,y]\), a principal ideal generated by a squarefree polynomial is radical, so
\[
k[x,y]/\bigl(y(y-cx-1)\bigr)
\]
is reduced for every \(c\).

The irreducible components are
\[
L_{1,c}=V(y),\qquad
L_{2,c}=V(y-cx-1).
\]
Each is isomorphic to \(\mathbb A^1_k\).  Thus neither the number nor the isomorphism type of the components changes with \(c\).

The incidence relation does change.  If \(c\ne0\), then the scheme-theoretic intersection has coordinate ring
\[
k[x,y]/(y,y-cx-1)
\cong
k[x]/(cx+1)
\cong
k.
\]
Hence the two lines meet in exactly one reduced point,
\[
(-c^{-1},0).
\]
If \(c=0\), the components are \(V(y)\) and \(V(y-1)\).  Their ideals are comaximal because
\[
1=y-(y-1),
\]
so the scheme-theoretic intersection is empty.  Equivalently,
\[
k[x,y]/(y(y-1))
\cong
k[x]\times k[x].
\]
Therefore the special fiber is the disjoint union of two parallel affine lines.

The complete degeneration audit is thus:
\[
\begin{array}{c|c|c|c}
c & \text{reduced?} & \text{components} & \text{incidence}\\ \hline
c\ne0 & \text{yes} & 2\text{ affine lines} & 1\text{ transverse point}\\
c=0 & \text{yes} & 2\text{ affine lines} & \text{disjoint}
\end{array}
\]
The degeneration changes how the components meet, while preserving reducedness and component count.
\end{solution}

\subsection*{Challenge 2: Arithmetic fiber classifier}
\begin{solution}
Let
\[
X=\Spec\mathbb Z[T]/(T^3-2)\longrightarrow\Spec\mathbb Z.
\]
At a prime \(p\),
\[
X_p=\Spec\mathbb F_p[T]/(T^3-2).
\]
We classify several primes by factoring the polynomial modulo \(p\).

For \(p=2\),
\[
T^3-2\equiv T^3,
\]
so
\[
X_2=\Spec\mathbb F_2[T]/(T^3).
\]
There is one point, with residue field \(\mathbb F_2\), but it has a nilpotent thickening of order \(3\).  This is a nonreduced fiber.

For \(p=3\),
\[
T^3-2\equiv T^3+1=(T+1)^3
\]
because \((T+1)^3=T^3+1\) in characteristic \(3\).  Hence
\[
X_3=\Spec\mathbb F_3[T]/((T+1)^3),
\]
again one nonreduced point of residue degree \(1\).

For \(p=5\), \(3^3\equiv2\pmod5\), and
\[
T^3-2=(T-3)(T^2+3T+4).
\]
The quadratic factor has discriminant
\[
3^2-4\cdot4=9-16\equiv3\pmod5,
\]
and \(3\) is not a square modulo \(5\).  Thus the quadratic is irreducible, giving
\[
X_5\cong
\Spec\mathbb F_5
\amalg
\Spec\mathbb F_{25}.
\]
The fiber is reduced with two points, of residue degrees \(1\) and \(2\).

For \(p=7\), the cubes modulo \(7\) are \(0,1,6\), so \(2\) is not a cube.  A cubic over a field is reducible only if it has a root, so \(T^3-2\) is irreducible over \(\mathbb F_7\).  Therefore
\[
X_7\cong\Spec\mathbb F_{7^3},
\]
one reduced point of residue degree \(3\).

For \(p=31\), \(4^3\equiv2\pmod{31}\), and because \(31\equiv1\pmod3\), the field contains the three cube roots of unity.  In fact
\[
T^3-2=(T-4)(T-7)(T-20)\quad\text{in }\mathbb F_{31}[T].
\]
Thus
\[
X_{31}\cong
\Spec\mathbb F_{31}
\amalg
\Spec\mathbb F_{31}
\amalg
\Spec\mathbb F_{31}.
\]

These examples display all the main scheme-theoretic data carried by factorization:
\[
\begin{array}{c|c|c|c}
p & \text{factorization type} & \text{residue degrees} & \text{reduced?}\\ \hline
2 & 1^3 & 1 & \text{no}\\
3 & 1^3 & 1 & \text{no}\\
5 & 1+2 & 1,2 & \text{yes}\\
7 & 3 & 3 & \text{yes}\\
31 & 1+1+1 & 1,1,1 & \text{yes}
\end{array}
\]
Repeated factors record ramification/nonreducedness, while degrees of distinct irreducible factors record residue degrees and splitting.
\end{solution}

\subsection*{Challenge 3: Geometric splitting versus nilpotence}
\begin{solution}
For the separable example, take
\[
K=\mathbb Q(\sqrt2).
\]
Then
\[
K\cong\mathbb Q[T]/(T^2-2).
\]
After extending scalars to an algebraic closure \(\overline{\mathbb Q}\),
\[
K\otimes_{\mathbb Q}\overline{\mathbb Q}
\cong
\overline{\mathbb Q}[T]/(T^2-2).
\]
The polynomial has two distinct roots \(\pm\sqrt2\), so the Chinese remainder theorem gives
\[
K\otimes_{\mathbb Q}\overline{\mathbb Q}
\cong
\overline{\mathbb Q}\times\overline{\mathbb Q}.
\]
Thus
\[
(\Spec K)_{\overline{\mathbb Q}}
\cong
\Spec\overline{\mathbb Q}\amalg\Spec\overline{\mathbb Q},
\]
two reduced geometric points.  The reason is separability: the minimal polynomial splits into distinct linear factors.

For the inseparable example, let
\[
k=\mathbb F_p(a^p),\qquad L=\mathbb F_p(a).
\]
Then
\[
L\cong k[T]/(T^p-a^p).
\]
After scalar extension to \(L\),
\[
L\otimes_kL
\cong
L[T]/(T^p-a^p)
=
L[T]/((T-a)^p).
\]
Writing \(\varepsilon=T-a\),
\[
L\otimes_kL\cong L[\varepsilon]/(\varepsilon^p).
\]
Its spectrum has one point, but \(\varepsilon\) is a nonzero nilpotent.

The contrast is therefore algebraically exact:
\[
\begin{array}{c|c|c}
\text{extension} & \text{specialized minimal polynomial} & \text{geometric effect}\\ \hline
\text{separable} & \text{distinct linear factors} & \text{several reduced points}\\
\text{purely inseparable} & \text{repeated }p\text{th power} & \text{nilpotent thickening}
\end{array}
\]
Geometric fibers distinguish splitting caused by separability from infinitesimal thickness caused by inseparability.
\end{solution}

\subsection*{Challenge 4: Fiber comparison under a tower}
\begin{solution}
Let
\[
S''\xrightarrow{h}S'\xrightarrow{g}S
\]
and let
\[
s''\mapsto s'\mapsto s.
\]
Set
\[
X'=X\times_SS',
\qquad
X''=X\times_SS''.
\]
The induced maps on residue fields form a tower
\[
\kappa(s)\longrightarrow\kappa(s')\longrightarrow\kappa(s'').
\]

Applying the fiber-after-base-change formula to \(S'\to S\) gives
\[
(X')_{s'}
\cong
X_s\times_{\kappa(s)}\kappa(s').
\]
Applying it directly to \(S''\to S\) gives
\[
(X'')_{s''}
\cong
X_s\times_{\kappa(s)}\kappa(s'').
\]
But \(X''\) is also the base change of \(X'\) from \(S'\) to \(S''\):
\[
X''\cong X'\times_{S'}S''.
\]
Therefore, applying the same formula over \(S'\),
\[
(X'')_{s''}
\cong
(X')_{s'}\times_{\kappa(s')}\kappa(s'').
\]
Substituting the first comparison,
\[
(X')_{s'}\times_{\kappa(s')}\kappa(s'')
\cong
\bigl(X_s\times_{\kappa(s)}\kappa(s')\bigr)
\times_{\kappa(s')}
\kappa(s'')
\cong
X_s\times_{\kappa(s)}\kappa(s'').
\]
Thus all routes give canonically isomorphic fibers:
\[
\boxed{
(X'')_{s''}
\cong
(X')_{s'}\times_{\kappa(s')}\kappa(s'')
\cong
X_s\times_{\kappa(s)}\kappa(s'')
}.
\]
The comparison maps are canonical because they come solely from associativity and uniqueness of fiber products; there is no coordinate choice involved.
\end{solution}

\subsection*{Challenge 5: Design a family with controlled special fiber}
\begin{solution}
Assume \(\operatorname{char}k\ne2\).

For the first family, take
\[
X=\Spec k[t,u]/(u^2-t^2)
\longrightarrow
\mathbb A^1_t.
\]
At \(t=a\ne0\),
\[
X_a=\Spec k[u]/((u-a)(u+a)).
\]
The two factors are distinct and comaximal, so
\[
X_a\cong\Spec(k\times k),
\]
two distinct reduced \(k\)-points.  At \(t=0\),
\[
X_0=\Spec k[u]/(u^2),
\]
one doubled nonreduced point.  Thus this family realizes the requested collision.

To obtain an empty special fiber while keeping two reduced points over every nonzero parameter, use an affine family supported over the principal open \(D(t)\):
\[
Y=\Spec k[t,t^{-1},u]/(u^2-1)
\longrightarrow
\mathbb A^1_t.
\]
For \(a\ne0\),
\[
Y_a
\cong
\Spec k[u]/(u^2-1)
=
\Spec k[u]/((u-1)(u+1))
\cong
\Spec(k\times k),
\]
again two distinct reduced points.

At \(t=0\), however,
\[
k[t,t^{-1},u]/(u^2-1)\otimes_{k[t]}k[t]/(t)=0,
\]
because the image of \(t\) is both invertible and zero.  Hence
\[
Y_0=\varnothing.
\]

The two constructions isolate two different mechanisms:
\[
\begin{array}{c|c}
\text{special fiber} & \text{algebraic mechanism}\\ \hline
\text{doubled point} & \text{two factors coalesce into }u^2\\
\text{empty} & \text{the parameter }t\text{ was inverted on the total space}
\end{array}
\]
Both are genuine scheme-theoretic changes in a family over the same affine base.
\end{solution}
